\documentclass{article}
\usepackage[utf8]{inputenc}

% TODO: Replace with your actual project title
\title{STAT 482 Proposal: [YOUR PROJECT TITLE HERE]}
% TODO: Replace with all team member names
\author{Team Member 1, Team Member 2, Team Member 3}
\date{Spring 2026}

\usepackage{natbib}
\usepackage{graphicx}


\begin{document}

\maketitle

%% ============================================================
%% SECTION 1: INTRODUCTION
%% This section should cover:
%%   - The data set and the question to be answered
%%   - The importance of the study
%%   - An introduction to the data set
%%
%% The entire proposal must be at least 300 words but no more
%% than 10 pages (including words, tables, and graphs).
%% ============================================================
\section{Introduction}

% TODO: Write 1--2 paragraphs introducing your topic area and why it matters.
% Explain the real-world context, the motivation for the study,
% and the specific research question(s) you aim to address.
[INTRODUCE YOUR TOPIC HERE. Describe the broader context of the problem you
are studying. Why is this topic important? What gap in knowledge or practical
need does your project address? Clearly state the research question(s) your
analysis will answer.]

% TODO: Introduce and describe the data set you will use.
% Where does it come from? How was it collected? What are the key variables?
[INTRODUCE YOUR DATA SET HERE. Describe where the data comes from (e.g., a
public repository, an API, a collaborator, etc.). Explain how the data were
collected, the time period covered, the number of observations, and the key
variables you will focus on. If the data set is publicly available, provide
a citation or URL.]

% --- Example of inline math (kept from the template for reference) ---
% You can include mathematical notation like this:
% $Y = f(\mathbf{s}, \mathbf{x}) + \epsilon, \quad \epsilon \stackrel{\mathrm{iid}}{\sim} \mathrm{N}(0, \sigma^{2})$

% --- Example of display-mode equation ---
% $$Y = f(\mathbf{s}, \mathbf{x}) + \epsilon, \quad \epsilon \stackrel{\mathrm{iid}}{\sim} \mathrm{N}(0, \sigma^{2})$$

% --- Example of a numbered equation ---
% \begin{equation}
%   Y = f(\mathbf{s}, \mathbf{x}) + \epsilon, \quad \epsilon \stackrel{\mathrm{iid}}{\sim} \mathrm{N}(0, \sigma^{2})
% \end{equation}

% --- Examples of citing references (using natbib) ---
% Parenthetical citation:  \citep{rigatti2017random}
% Textual citation:        \cite{rigatti2017random}
% See ProposalRef.bib for adding your own references.

% TODO: Provide a brief literature review or background.
% Identify relevant prior work, methods, or studies that motivate your approach.
[PROVIDE BACKGROUND / LITERATURE REVIEW HERE. Discuss any relevant prior work
or existing studies. What methods have been used before? What are the
limitations of previous approaches? How does your project build upon or differ
from this prior work? Cite references as appropriate using
\texttt{\textbackslash citep\{\}} and \texttt{\textbackslash cite\{\}}.]

% TODO: Clearly state the objectives of your project.
[STATE YOUR PROJECT OBJECTIVES HERE. Summarize the specific goals of your
analysis in a bulleted or numbered list.]

\begin{itemize}
    \item[1.] [OBJECTIVE 1: e.g., Build a predictive model for \ldots]
    \begin{itemize}
        \item [Sub-goal 1a]
        \item [Sub-goal 1b]
    \end{itemize}
    \item[2.] [OBJECTIVE 2: e.g., Compare the performance of \ldots]
\end{itemize}


%% ============================================================
%% SECTION 2: RESEARCH STRATEGY
%% This section should cover:
%%   - A detailed description of the data
%%   - The statistical analysis plan
%% ============================================================
\section{Research Strategy}\label{sec:research}

\subsection{Data}

% TODO: Provide a detailed description of your data set.
% Include information such as: source, sample size, variables, data types,
% any preprocessing or cleaning needed, and any known issues or limitations.
[DESCRIBE YOUR DATA IN DETAIL HERE. Discuss the source, how it was collected,
the sample size, the variables of interest (response variable and predictors),
data types, any missing data, and any preprocessing steps you plan to take.
If you have preliminary exploratory data analysis (EDA), include summary
statistics or initial visualizations.]

% --- Example of including a figure (kept from the template for reference) ---
% Replace Trace.pdf with your own figure file.
% \begin{figure}[h!]
% \centering
% \includegraphics[width=0.45\textwidth]{your_figure.pdf}
% \caption{[DESCRIBE YOUR FIGURE HERE]}
% \label{fig:your_label}
% \end{figure}

% TODO: If you have a figure to include, uncomment the block above and
% replace the filename and caption. You can reference it in the text
% with \ref{fig:your_label}.

\subsection{Planned Analysis}

% TODO: Describe the statistical methods and analysis plan you intend to use.
% Be specific about models, techniques, software, and evaluation criteria.
[DESCRIBE YOUR STATISTICAL ANALYSIS PLAN HERE. What statistical methods do
you plan to use (e.g., linear regression, logistic regression, time series
analysis, random forests, clustering, etc.)? Why are these methods appropriate
for your data and research questions? Describe any model assumptions, how you
will check them, and how you will evaluate model performance (e.g.,
cross-validation, AIC/BIC, residual diagnostics, etc.).]

% TODO: Describe any additional analyses or extensions you plan to explore.
[DESCRIBE ANY ADDITIONAL ANALYSES HERE. For example, do you plan to compare
multiple models? Perform variable selection? Conduct sensitivity analyses?
Discuss how the results will be interpreted and presented.]


%% ============================================================
%% BIBLIOGRAPHY
%% Add your references to ProposalRef.bib in BibTeX format.
%% The template uses the chicago bibliography style.
%% ============================================================
\bibliographystyle{chicago}
\bibliography{ProposalRef.bib}

\end{document}
